\BeleskeID{01-s019}
% element: 01-s019:e01
% element: 01-s019:e02
% element: 01-s019:e03

% element: 01-s019:d01
Pri računanju NI funkcije najpre odredimo proizvod ulaza, pa negiramo ceo proizvod. Negacioni kružić na izlazu predstavlja taj poslednji korak. Ne treba ga tumačiti kao negaciju samo jednog ulaza.

Ako je jedan ulaz 0, proizvod je 0, pa je izlaz NI kola 1. Ako je jedan ulaz stalno 1, preostali ulaz se negira: \(\overline{x\cdot1}=\bar{x}\). To je koristan način da se iz iste funkcije prepozna invertorsko ponašanje.
